\scanpage{49}\mspage{48}

Ces relations sont symétriques, \unsure{.} $\mathrel{\rlap{$>$}<}$ est réflexive (tautologique), $\underset{M}{\parallel}$ antiréflexive (utiliser le fait que $\underset{\Lambda}{\parallel}$ l'est, et $M\Lambda2$, impliquant que $\mathrm{omb}_{\Lambda}(X)\ne\varnothing$). Je dis

\uline{Lemme 1} \unsure{Soient|Soit} $X,Y\in M$.
\[
\begin{aligned}
\text{a)}\quad X\underset{M}{\parallel}Y &\Longleftrightarrow X\mathrel{\rlap{$>$}<} Y,\quad\text{et}\quad \underbrace{\widetilde X\cap\widetilde Y}_{L}=\varnothing\\
\text{b)}\quad\text{Si } X,Y\in\Lambda,\quad&\text{alors}\quad X\underset{M}{\parallel}Y\Longleftrightarrow X\underset{\Lambda}{\parallel}Y
\end{aligned}
\]
(on omettra donc par la suite les indices $M,\Lambda$ dans $\parallel$)

Prouvons a). \deleted{\unsure{On a} $\Longrightarrow$ : $X\mathrel{\rlap{$>$}<} Y$ \unsure{trivial sur} définition, $\widetilde X\cap\widetilde Y=\varnothing$}

Prouvons $\Longrightarrow$ i.e. supposons $X\underset{M}{\parallel}Y$, prouvons :
\[
\begin{aligned}
1^{\circ})\quad &X\mathrel{\rlap{$>$}<} Y:\ \text{tautologique}\\
2^{\circ})\quad &\widetilde X\cap\widetilde Y=\varnothing,\quad\text{si on avait }Z\in\widetilde X\cap\widetilde Y,\ \text{soit}
\end{aligned}
\]
$Z'\in\mathrm{omb}_{\Lambda}(Z)$, alors $Z'\in\mathrm{omb}_{\Lambda}(X)\cap\mathrm{omb}_{\Lambda}(Y)$ \deleted{\unsure{donc par}.} $X\parallel Y$ on doit avoir $Z'\parallel Z'$, absurde (antiréflexivité).

Prouvons $\Longleftarrow$ i.e. supposons $X\mathrel{\rlap{$>$}<} Y$, $\widetilde X\cap\widetilde Y=\varnothing$. et prouvons $X\parallel Y$. \inserted{Mais comme \deleted{par déf.}} \deleted{Soit $X'\in\mathrm{omb}_{\Lambda}(X)$, $Y'\in\mathrm{omb}_{\Lambda}(Y)$.} \deleted{, de $X\mathrel{\rlap{$>$}<} Y$, \unsure{donc le}} \unsure{on|a} $L=\widetilde X\cap\widetilde Y=\varnothing$, \deleted{\unsure{alr|alors}} donc $\mathrm{omb}_{\Lambda}(L)=\varnothing$, \deleted{\unsure{elle se réduit à|il en résulte} : $\forall X'\in\mathrm{omb}_{\Lambda}(X)$} la définition de $X\mathrel{\rlap{$>$}<} Y$ montre que ça \unsure{signifie|signifi:} $X\parallel Y$.

\uline{Prouvons} b Résulte tautologiquement de $M\Lambda4$

\uline{Terminologie}\unsure{.} Soi\deleted{\unsure{t}}\unsure{ent|t} $P,Q$ deux parties de $M$. On pose
\[
\begin{aligned}
P\parallel Q&\overset{\mathrm{def}}{\Longleftrightarrow}\quad\forall X\unsure{\mathpunct{:}}\in P,\ \forall Y\in Q,\ \text{\inserted{on a}}\quad X\parallel Y\\
&\Longleftrightarrow\quad\forall X'\in\mathrm{omb}_{\Lambda}P,\ Y'\in\mathrm{omb}_{\Lambda}(Q),\ \text{\inserted{on a}}\quad X'\parallel Y'
\end{aligned}
\]
(où $\mathrm{omb}_{\Lambda}(P)=\bigcup_{X\in P}\mathrm{omb}_{\Lambda}(X)$, et \unsure{item|idem} pour $\mathrm{omb}_{\Lambda}(Q)$)
